\section{A common interface for octagon capacities}\label{sec:octagons}
For the octagon arguments, coordinates are \emph{cap first, cup second}.
This is the transpose of the heptagon certificate convention, not a
change in the parameter order of $M_n(a,b)$.

\begin{definition}[Admissible labels and states]\label{def:admissible}
Let $P$ be a planar point set in general position. Assign a label
$d(u,v)=(p(u,v),q(u,v))\in\mathbb Z_{>0}^2$ to every ordered pair
$u<v$, and a finite down-set $I_v\subset\mathbb Z_{>0}^2$ to every
point. The assignment is admissible if
\begin{align}
 p(u,v)&<p(v,w)&&\text{for each cap triple }u<v<w,\label{eq:growthcap}\\
 q(u,v)&<q(v,w)&&\text{for each cup triple }u<v<w,\label{eq:growthcup}\\
 d(u,v)&\in I_v\setminus I_u&&\text{for each }u<v.\label{eq:admissible}
\end{align}
A family of states has \emph{admissible capacity} $b$ if every
octagon-free set with any such assignment occupies at most $b$ of
its states.
\end{definition}

The states are distinct, and strict inclusion forces the same owner
order: the opposite order would make \eqref{eq:admissible} an empty
directed difference.

Ordinary actual incoming ranks, with their incoming down-sets, provide
an admissible assignment by \cref{lem:difference}. More generally,
labels satisfying \eqref{eq:growthcap}--\eqref{eq:growthcup} provide one
by taking $I_v=\down\{d(u,v):u<v\}$. The same domination proof applies.
The states in an admissible assignment need not equal these generated
down-sets. Restricting to a subset of points and \emph{retaining} its
labels and assigned states preserves admissibility. This is the
interface needed for transport; a bound using actual attainment cannot
automatically be transported in this way.

\begin{lemma}[Forced signs]\label{lem:forced}
Suppose the chronological owners of states $A,B,C$ exist, and put
$\Delta_1=B\setminus A$, $\Delta_2=C\setminus B$. If
\[
 \max_{z\in\Delta_2}z_2\le\min_{z\in\Delta_1}z_2,
\]
their triple is a cap. Interchanging coordinates forces a cup.
Both directed differences must be nonempty.
\end{lemma}
\begin{proof}
The two consecutive pair labels belong to $\Delta_1$ and $\Delta_2$.
A cup would strictly increase the second coordinate, contradicting the
displayed inequality. The other assertion is dual.
\end{proof}

\subsection{Hereditary convex families}
A cyclic state list is certified by checking every triple in its cyclic
order under all six chronological permutations. Empty directed pair
differences exclude an order. Simultaneous forcing of both signs also
excludes an order. Every other order must have the sign required by the
cyclic orientation, with permutation parity accounted for.

\Needspace{7\baselineskip}
\begin{lemma}[Hereditary convexity]\label{lem:cycle}
Every occupied subset of a certified cyclic family is in convex
position. Consequently such a family $K$ has admissible capacity seven.
\end{lemma}
\begin{proof}
For at least three occupied states, every triple of their owners is
counterclockwise in the inherited cyclic order. For each directed edge
between consecutive selected vertices, every other selected vertex is
strictly on its left. Each edge is supporting, so all selected vertices
are hull vertices. Eight occupied states would give an octagon.
\end{proof}

This test is hereditary because a triple's test uses only those three
states. A proof for a whole path with fixed endpoints would not, by
itself, justify a capacity after some endpoints are omitted.

\section{The two octagon omission certificates}\label{sec:319}
Write $h_A=1$ when state $A$ is absent. A valid occupied capacity $b_K$
on a family $K$ gives the omission inequality
\[
 \sum_{A\in K}h_A\ge |K|-b_K.
\]
We first give the arithmetic rule used in both finite proofs.

\Needspace{13\baselineskip}
\begin{lemma}[Exact weighted contradiction]\label{lem:weighted}
Suppose omission rows are $\sum_{v\in K}h_v\ge a_K$. At a branch node,
let $O$ be the states fixed missing, $Z$ the states fixed occupied,
$V$ the unassigned states, and $k$ the remaining hole budget. For
nonnegative integer weights $w_K$ and an integer $\Lambda>0$, put
\begin{align*}
 r_K&=a_K-|K\cap O|,& W&=\sum_Kw_Kr_K,\\
 \ell(v)&=\sum_{K\ni v}w_K,&
 P&=\sum_{v\in V}\max(0,\ell(v)-\Lambda).
\end{align*}
Every completion using at most $k$ further holes satisfies $W\le \Lambda k+P$.
Thus $W>\Lambda k+P$ certifies the node as impossible.
\end{lemma}
\begin{proof}
Adding the weighted residual inequalities and using binary values gives
\[
 W\le\sum_{v\in V}\ell(v)h_v
 \le\Lambda\sum_{v\in V}h_v+\sum_{v\in V}\max(0,\ell(v)-\Lambda)
 \le \Lambda k+P.
\]
All operations are exact integer arithmetic; $P$ accounts for excess load.
\end{proof}

\subsection{The balanced bound}
Ordinary cap/cup incoming states inject a balanced avoiding set into
the 252 down-sets of $[5]^2$. The supplied catalogue contains 2,685
families: 1,584 of size eight, 780 of size nine, and 321 of size ten.
Every family has a certified cyclic order and hence capacity seven by
\cref{lem:cycle}. 

For example, the following ten-state cyclic list requires at least
three omissions:
\begin{align*}
 &(3,2,1,0,0),\ (3,2,1,1,1),\ (3,2,2,1,1),\ (3,3,2,1,1),\\
 &(4,3,2,1,1),\ (5,3,2,1,1),\ (5,3,2,1,0),\ (5,3,2,0,0),\\
 &(5,3,1,0,0),\ (5,2,1,0,0).
\end{align*}
The check covers all 120 triples and their chronological permutations,
as illustrated in Figure~\ref{fig:states}. The independent checker
rejects 472 of the 720 orders and forces the required sign in all
248 others.

\begin{figure}[ht]
\centering
\newcommand{\ferrerspic}[1]{\begin{tikzpicture}[x=.16cm,y=.16cm,baseline=0cm]
 \path[use as bounding box] (0,0) rectangle (5,5);
 \foreach \h [count=\i from 0] in {#1}{
  \ifnum\h>0\foreach \j in {1,...,\h}{
   \draw[fill=black!15,line width=.25pt] (\i,\j-1) rectangle +(1,1);}\fi}
 \end{tikzpicture}}
\begin{tikzpicture}[>=Stealth, every node/.style={inner sep=3pt}]
 \node (a1) at (0,1.9) {\ferrerspic{3,2,1,0,0}};
 \node (a2) at (1.65,1.9) {\ferrerspic{3,2,1,1,1}};
 \node (a3) at (3.3,1.9) {\ferrerspic{3,2,2,1,1}};
 \node (a4) at (4.95,1.9) {\ferrerspic{3,3,2,1,1}};
 \node (a5) at (6.6,1.9) {\ferrerspic{4,3,2,1,1}};
 \node (a6) at (6.6,0) {\ferrerspic{5,3,2,1,1}};
 \node (a7) at (4.95,0) {\ferrerspic{5,3,2,1,0}};
 \node (a8) at (3.3,0) {\ferrerspic{5,3,2,0,0}};
 \node (a9) at (1.65,0) {\ferrerspic{5,3,1,0,0}};
 \node (a10) at (0,0) {\ferrerspic{5,2,1,0,0}};
 \foreach \i/\j in {1/2,2/3,3/4,4/5,5/6,6/7,7/8,8/9,9/10,10/1}
   \draw[->] (a\i)--(a\j);
 \foreach \i in {1,...,5}\node[font=\scriptsize,above=1pt of a\i] {\i};
 \foreach \i in {6,...,10}\node[font=\scriptsize,below=1pt of a\i] {\i};
\end{tikzpicture}
\caption{The ten Ferrers states in their certified cyclic order.
The arrows specify the order of the certificate, not the chronological
order of points. At most seven states can be occupied.}
\label{fig:states}
\end{figure}

\begin{proposition}[Balanced octagon bound]\label{prop:193}
$C_8=M_8(7,7)\le193$.
\end{proposition}
\begin{proof}
The complete binary certificate assumes at most 58 missing states.
It has 1,590 branches and 1,591 weighted leaves, with maximum depth
14. Each split assigns one previously unassigned Boolean variable
and includes both children. Every node is reached exactly once, and
all leaves satisfy \cref{lem:weighted} in their actual branch context.
There are 306,922 positive weight terms and 119 leaves with $P>0$.
The minimum strict margin is 21.

For the smallest-margin leaf, the remaining budget is $k=52$,
$\Lambda=10^6$, $W=52{,}000{,}021$, and $P=0$. Thus a completion would
require $52{,}000{,}021\le52{,}000{,}000$. The other 1,590 leaves
are checked in the same exact fashion. Every hypothetical hole set
with size at most 58 follows a path to a contradiction. At least 59
states are absent, leaving at most $252-59=193$ points.
\end{proof}

\subsection{A two-sided residual state space}
For a set $R$ counted by $D_8$, let $C^+(u,v)+1$ be the longest cap
starting at the exact pair $(u,v)$ and $Q^-(u,v)+1$ the longest cup
ending there, both in the entire fixed $R$. Define
\begin{equation}\label{eq:twosided}
 d(u,v)=(5-C^+(u,v),Q^-(u,v)),\qquad
 \Psi(v)=\down\{d(u,v):u<v\}.
\end{equation}
These coordinates lie in $[4]\times[6]$. A cap triple $u<v<w$ gives
$C^+(u,v)\ge C^+(v,w)+1$, so the first coordinate strictly increases;
a cup triple strictly increases the second. Thus these labels and
states are admissible and injective. Future-cap statistics have
precedent in Baek's work~\cite[Theorem~3.2 and Definition~3.4]{Baek};
the useful specialization here is to the residual forbidden junction.

\begin{proposition}[Residual state restriction]\label{prop:154}
Every residual state belongs to the 154-element universe
\[
 \mathcal U=\left\{h\in\mathbb Z^4:\begin{array}{l}
 6\ge h_1\ge h_2\ge h_3\ge h_4\ge0,\\
 h_1<6\text{ or }h_1=h_2=6
 \end{array}\right\}.
\]
\end{proposition}
\begin{proof}
No actual pair label is $(1,6)$. Such a label would exhibit a five-cap
starting at $(u,v)$ and a seven-cup ending at $(u,v)$. Removing $v$
from the latter gives a six-cup ending at $u$, where the five-cap starts.
This is the forbidden residual junction.

Every maximal cell of an actual incoming down-set is attained by
some incoming pair. Hence $(1,6)$ cannot be maximal. Among the
$\binom{10}{4}=210$ down-sets of $[4]\times[6]$, exactly those with
$h_1=6$ and $h_2\le5$ have that maximal cell. There are
$\binom83=56$ such shapes, leaving $210-56=154$.
\end{proof}

A state containing $(1,6)$ remains allowed when a larger cell
dominates it. An exact
154-point control realizes all the allowed states while satisfying
the chain and junction restrictions; it contains a checked octagon.
Thus further omissions genuinely require polygon avoidance.

\begin{proposition}[General residual baseline]\label{prop:baseline}
For $n\ge6$,
\[
 D_n\le\binom{2n-6}{n-4}-\binom{2n-8}{n-5}.
\]
Consequently, the cup--cap estimate for $C_n$ and
\cref{thm:transfer} recover the finite bound \eqref{eq:MV}.
\end{proposition}
\begin{proof}
Use the labels $(n-3-C^+(u,v),Q^-(u,v))$ in
$[n-4]\times[n-2]$, with both ranks computed in the residual set.
The growth and injection argument is unchanged. A label $(1,n-2)$
would give an $(n-3)$-cap starting at $u$ and an $(n-1)$-cup ending
at the pair $(u,v)$, hence the forbidden junction at $u$.
So $(1,n-2)$ cannot be a maximal cell. The excluded down-sets have
first column $n-2$ and remaining $n-5$ columns of height at most
$n-3$, giving $\binom{2n-8}{n-5}$ exclusions.
Finally, Pascal's identity and binomial symmetry give
\begin{align*}
 &\binom{2n-6}{n-3}+\binom{2n-6}{n-4}-\binom{2n-8}{n-5}+2\\
 &\hspace{25mm}=\binom{2n-5}{n-2}-\binom{2n-8}{n-3}+2.
\end{align*}
This is a derivation of the known baseline, not a new bound for general $n$.
\end{proof}

\subsection{Transporting four-cap capacities}\label{sec:compression}
The following lemma supplies numerical information from Baek's theorem
without assuming that transported labels are actual longest ranks.

\begin{lemma}[Horizontal block compression]\label{lem:compression}
For $\rho=1,\ldots,s$, let nonnegative integer thresholds
$f_\rho(0)<\cdots<f_\rho(r)$ lie in a larger rank rectangle and satisfy
\[
 f_{\rho+1}(i)\le f_\rho(i),\qquad
 f_{\rho+1}(i)\le f_\rho(i-1)+1\quad(1\le i\le r),
\]
where the first inequality also includes $i=0$. Map a source ideal
with row widths $w_\rho$ to the ideal with row widths $f_\rho(w_\rho)$.
In any admissible assignment, the owners of states in this image
admit labels in $[r]\times[s]$ satisfying strict growth. Consequently
they contain neither an $(r+2)$-cap nor an $(s+2)$-cup.
\end{lemma}
\begin{proof}
Threshold monotonicity makes every image a down-set, and strictness
makes the state map injective. Directed differences are precisely the
unions of intervals belonging to the corresponding source difference.
Compress a pair label $(R,\rho)$ to $(i,\rho)$ when
$f_\rho(i-1)<R\le f_\rho(i)$. Such an interval exists uniquely because
the label is in a directed image difference.

Cup growth is unchanged. At a cap triple, suppose the outgoing
compressed cell $(j,\sigma)$ had $j\le i$, where $(i,\rho)$ is incoming.
The middle source ideal contains the incoming cell and excludes the
outgoing one, so its down-set property forces $\sigma>\rho$. Therefore
\[
 R_{\rm out}\le f_\sigma(j)\le f_{\rho+1}(i)
 \le f_\rho(i-1)+1\le R_{\rm in},
\]
contrary to cap growth. Thus the compressed first coordinate increases
on every cap triple. Bounded strictly increasing ranks exclude the
stated chains.
\end{proof}

With $r=2$ and $s=6$, the source has 28 states. The image owners have
no four-cap and inherit octagon avoidance. Baek's
theorem~\cite[Theorem~1.6]{Baek} says that 23 points force a four-cap
or an octagon. Each image therefore has occupied capacity 22, and
requires six omissions. The residual catalogue contains 23 explicit
maps into $\mathcal U$. The finite validation checks the threshold conditions, images,
and directed differences for all 23 maps.

\subsection{The residual omission theorem}
The final formal catalogue has 1,187 omission inequalities:
1,164 hereditary convex families, each with demand $|K|-7$, and
23 four-cap block images, each with demand six. The original discovery
archive split the first group into 973 transported families and 191
cycles. In the formalization, every one of the 973 transported supports
also receives an explicit cyclic certificate. Thus all 1,164 rows use
\cref{lem:cycle}; their historical transport provenance is no longer
an additional proof dependency. The remaining 23 rows use
\cref{lem:compression} and the port of Baek's theorem.

\begin{proposition}[Residual octagon bound]\label{prop:124}
$D_8\le124$.
\end{proposition}
\begin{proof}
Assume at most 29 of the 154 states are missing. The complete tree has
7,314 binary branches and 7,315 weighted leaves, maximum depth 21.
Every split includes both values of one previously unassigned state,
and every node is reached exactly once. Each leaf satisfies
\cref{lem:weighted} with its recorded context. There are 662,561
positive weight terms; the minimum strict integer margin is one.
Thus no Boolean vector with at most 29 holes satisfies all the valid
inequalities. At least 30 states are absent and $|R|\le154-30=124$.
\end{proof}

Finally, \cref{thm:transfer,prop:193,prop:124} give
\[
 Q_8\le193+124=317,\qquad
 \boxed{\ES(8)\le193+124+2=319}.
\]
The classical lower construction gives $\ES(8)\ge65$~\cite{ES61}.
